\section{Background}

\subsection{What the system is}

Ariadne answers questions about attacker behaviour by querying two live graph
databases and joining the results. MITRE ATT\&CK sits in Neo4j and is queried in
Cypher. MITRE CAPEC sits in a TinkerPop Gremlin Server and is queried in
Gremlin. The two are joined on shared ATT\&CK technique identifiers. Every
sentence in an answer carries a citation to a retrieved record and to the query
that produced it.

The pipeline plans a question into sub-questions, routes each to a source,
retrieves, verifies each claim against the retrieved records, and composes an
answer. Almost all of that is deterministic. Query construction fills a fixed
library of parameterised, read-only templates rather than generating query text.

\subsection{Why the project exists in this form}

The design commitment is unfashionable. A reviewer who expects an agent sees
hand-written templates and pattern-matched intent classification. The project's
own external audit concluded that the way to defend this is to stop selling the
mechanism and start selling the measurement \citep{audit2026}.

That reframing has support in the literature. Writing graph queries from English
is measurably hard. Zero-shot execution accuracy on an unseen production graph
averages 0.036 across 11 models \citep{pipecypher}. Zero-shot execution accuracy
on ISO-GQL is at or below 4\% \citep{text2gql}. A peer-reviewed cross-dialect
benchmark reports 35.40\% average execution accuracy \citep{gqlbench}.

Choosing templates over generation is not new either. It is a named family in
knowledge-base question answering since 2015 \citep{aqqu, tqa}, catalogued as
such in the field's own survey \citep{kbqasurvey}. The modern analogue
constrains decoding rather than eliminating it, and it works \citep{picard}.

\subsection{The gap}

Two things were missing when this project started, and the audit verified both
across seven search formulations.

First, nobody had run a controlled head-to-head of a query-template library
against free-form query generation for graph query languages. The comparison is
asserted in a relational setting \citep{finai} and never measured for Cypher or
Gremlin.

Second, nobody had reported federated question answering across security
knowledge graphs held in separate engines with different query languages. The
closest prior work, AgCyRAG, builds that architecture in the same domain and
commits to predefined rather than generated queries \citep{agcyrag}. It reports
no quantitative evaluation at all. It scoops the design and cannot scoop the
result.

\subsection{The questions this record set out to answer}

\begin{enumerate}
  \item Does routing to validated templates beat letting a language model write
        the query?
  \item Is the pipeline's near-perfect score a real capability, or an artifact
        of a question set written alongside the code that answers it?
  \item Does querying two engines separately beat merging both sources into one
        graph?
\end{enumerate}

\subsection{Scope}

This report covers work recorded between 2026-07-20 and 2026-08-20. It covers
one domain, two catalogues, and one served language model. It compares Ariadne
against versions of Ariadne. No external system has been implemented or run, and
question 3 is unanswered.
